\section{Introduction}

A coding agent starts each session knowing nothing it learned in the last one. Native memory is a set of files loaded whole: Claude Code loads its \texttt{CLAUDE.md} files and the first 200 lines or 25\,KB of a memory index, with no ranking by the task at hand \cite{claudecode_memory}. Memory products add retrieval, but they select by similarity top-$k$ or leave the choice to the agent, and they spend LLM calls on every write. The published evidence does not show that this is good enough. On LoCoMo, Mem0's own table has a full-context baseline ahead of every memory system tested \cite{chhikara2025mem0}. Retrieval method moves accuracy far more than write strategy, and raw chunks with no LLM calls match fact extraction \cite{yuan2026diagnosing}. On coding tasks, unfiltered experience gives little or negative benefit \cite{zhu2026swecontextbench}, and eleven of twelve pairings of solvers with memory systems fail to beat memory-off \cite{vibemembench2026}.

These findings point at the read path, and at three defects in it. Selection ignores what an item costs in tokens. It ignores what the other selected items already say. And it ignores whether an item is still true, which is the common failure in a codebase, where a function signature read last month may have changed since.

This paper treats memory as an online budgeted selection problem and makes four contributions.
\begin{enumerate}
\item A framework in which the surveyed systems are special cases, with elementary worst-case arguments for what each one gives up (Section~\ref{sec:framework}).
\item A method that separates \emph{validity} from \emph{usefulness}: budgeted monotone submodular selection with a per-query optimality certificate, a validity model estimated from the store's own version history, and a linear policy learned from execution feedback (Section~\ref{sec:method}). It runs locally and needs no LLM.
\item A regret decomposition for items that expire (Theorem~\ref{thm:regret}). Stationary linear submodular bandit results exist \cite{yue2011linear,yu2016linear,takemori2020submodular}; they exclude this case. Here non-stationarity costs the accumulated validity error, while a validity-blind learner has linear regret (Proposition~\ref{prop:blind}). The other theory results are a rate--distortion bound for reversible consolidation and a stability result for store size.
\item RepoMem, a coding-memory benchmark built from the real histories of eight Python repositories, in which facts change for real and every label comes from the code by static analysis. We evaluate on it, on LongMemEval-S and on LoCoMo against a dozen reading baselines and Mem0's verbatim mode, with ablations, confidence intervals and the cases where the method loses.
\end{enumerate}

\paragraph{Scope of the evidence.} All results here are \emph{retrieval-level}: evidence recall under a token budget, tokens injected, staleness errors, latency and store size. No LLM was available for these experiments, so we report no answer accuracy and no end-to-end task success. Recall under a budget is a necessary condition for those outcomes, not a measurement of them. A controlled study found that context strategy moved correctness on coding tasks by no more than 10 to 15 points \cite{khatri2026contextfiles}, so large end-to-end gains should not be assumed.
